\subsection{Мотивация функции потерь MSE}

Мы начнём с использования линейных моделей для решения задачи регрессии.
Простейшим примером постановки задачи линейной регрессии является метод наименьших квадратов.

Пусть у нас задан датасет $(X,~y)$, где $y = (y_i)_{i=1}^N \in \mathbb{R}^N$ -- вектор значений целевой переменной,
а $X = (x_i)_{i=1}^N \in \mathbb{R}^{N \times D},~x_i \in  \mathbb{R}^D$ -- матрица объекты-признаки, в которой $i$-я строка --
это вектор признаков $i$-го объекта выборки.
Мы хотим моделировать зависимость $y_i$ от $x_i$ как линейную функцию со свободным членом.
Общий вид такой функции из $\mathbb{R}^D$ в $\mathbb{R}$ выглядит следующим образом:

\[ f_w(x_i) = (w,~x_i) + w_0 \]

Свободный член $w_0$ часто опускают, потому что такого же результата можно добиться, добавив ко всем $x_i$ признак,
тождественно равный единице; тогда роль свободного члена будет играть соответствующий ему вес:

\[(x_{i_1} \dots x_{i_D}) \dot (w_1 \dots w_D)^T + w_ 0 =(1~x_{i_1} \dots x_{i_D}) \dot (w_0~w_1 \dots w_D)^T\]

Поскольку это сильно упрощает запись, в дальнейшем мы будем считать, что это уже сделано и зависимость имеет вид
просто $f_w(x_i) = (w,~x_i)$.
Мы хотим, чтобы на нашем датасете (то есть на парах $(x_i,~y_i)$ из обучающей выборки) функция $f_w$ как можно
лучше приближала нашу зависимость

Для того чтобы чётко сформулировать задачу, нам осталось только одно: на математическом языке выразить желание
«приблизить $f_w(x)$ к $y$».
Говоря простым языком, мы должны научиться измерять качество модели и минимизировать её ошибку, меняя обучаемые
параметры.
В нашем примере обучаемые параметры -- это веса $w$.
Функция, оценивающая то, как часто модель ошибается, традиционно называется функцией потерь, функционалом
качества или просто loss (loss function).
Важно, чтобы её было легко оптимизировать: гладкая функция потерь -- это хорошо, а кусочно постоянная -- ужасно.

Функции потерь бывают разными.
От их выбора зависит то, насколько задачу в дальнейшем легко решать, и то, в каком смысле у нас получится
приблизить предсказание модели к целевым значениям.
Интуитивно понятно, что для нашей текущей задачи нам нужно взять вектор $y$ и вектор предсказаний модели и
как-то сравнить, насколько они похожи.
Так как эти векторы «живут» в одном векторном пространстве, расстояние между ними вполне может быть функцией потерь.
Более того, положительная непрерывная функция от этого расстояния тоже подойдёт в качестве функции потерь.
При этом способов задать расстояние между векторами тоже довольно много.
От всего этого разнообразия глаза разбегаются, но мы обязательно поговорим про это позже.
Сейчас давайте в качестве loss возьмём квадрат $L_2$-нормы вектора разницы предсказаний модели и $y$.
Во-первых, как мы увидим дальше, так задачу будет нетрудно решить, а во-вторых, у этого loss есть ещё несколько
дополнительных свойств

\begin{itemize}

    \item $L_2$-норма разницы -- это евклидово расстояние $\abs{y - f_w(x)}^2$ между вектором таргетов и
    вектором ответов модели, то есть мы их приближаем в смысле самого простого и понятного «расстояния».

    \item Как мы увидим в разделе про вероятностные модели, с точки зрения статистики это соответствует гипотезе
    о том, что наши данные состоят из линейного «сигнала» и нормально распределенного «шума».

\end{itemize}

Так вот, наша функция потерь выглядит так:

\[ L(f,X,y) = \abs{y - f_w(x)}^2_2 = \abs{y - X w}^2_2 = \sum_{i=1}^N (y_i - (x_i,w))^2 \]

Такой функционал ошибки не очень хорош для сравнения поведения моделей на выборках разного размера.
Представьте, что вы хотите понять, насколько качество модели на тестовой выборке из 2500 объектов хуже, чем на
обучающей из 5000 объектов.
Вы измерили $L_2$-норму ошибки и получили в одном случае 300, а в другом 500.
Эти числа не очень интерпретируемы.
Гораздо лучше посмотреть на среднеквадратичное отклонение

\[ L(f,X,y) = \frac{1}{N} \sum_{i=1}^N (y_i - (x_i,w))^2 \]

По этой метрике на тестовой выборке получаем 0,12, а на обучающей 0,10.

Функция потерь $\frac{1}{N} \sum_{i=1}^N (y_i - (x_i,w))^2$ называется Mean Squared Error, MSE или
среднеквадратическим отклонением.
Разница с $L_2$-нормой чисто косметическая, на алгоритм решения задачи она не влияет:

\[ MSE(f,X,y) = \frac{1}{N} \abs{y - X w}^2_2 \]

В самом широком смысле, функции работают с объектами множеств: берут какой-то входящий объект из одного множества и
выдают на выходе соответствующий ему объект из другого.
Если мы имеем дело с отображением, которое на вход принимает функции, а на выходе выдаёт число, то такое отображение
называют функционалом.
Если вы посмотрите на нашу функцию потерь, то увидите, что это именно функционал: для каждой конкретной линейной
функции, которую задают веса $w_i$, мы получаем число, оценивающее, насколько точно эта функция приближает наши
значения $y$.
Чем меньше это число, тем точнее наше решение, значит для того, чтобы найти лучшую модель, этот функционал нам надо
минимизировать по $w$:

\[ \abs{y - X w}^2_2 \rightarrow \min_w \]

\subsubsection{Доказательство для гауссовского случая}

\underline{Гауссово предположение}
Предположим, что ошибки $\varepsilon_i$:
\begin{enumerate}
    \item Независимы и одинаково распределены (i.i.d.).
    \item Следуют нормальному распределению: $\varepsilon_i \sim \mathcal{N}(0, \sigma^2)$.
    \item Имеют нулевое математическое ожидание: $\mathbb{E}[\varepsilon_i] = 0$.
    \item Имеют постоянную дисперсию: $\text{Var}(\varepsilon_i) = \sigma^2 < \infty$.
\end{enumerate}

Тогда условное распределение $y_i$ при заданном $\bm{x}_i$ также нормальное:
\[ y_i \mid \bm{x}_i \sim \mathcal{N}(\bm{x}_i^\top \bm{\beta}, \sigma^2) \]

\underline{Функция правдоподобия}.
Функция правдоподобия для всех наблюдений при условии независимости:
\begin{align*}
    L(\bm{\beta}, \sigma^2) &= p(\bm{y} \mid \bm{X}; \bm{\beta}, \sigma^2) \\
    &= \prod_{i=1}^n p(y_i \mid \bm{x}_i; \bm{\beta}, \sigma^2) \\
    &= \prod_{i=1}^n \frac{1}{\sqrt{2\pi\sigma^2}} \exp\left(-\frac{(y_i - \bm{x}_i^\top \bm{\beta})^2}{2\sigma^2}\right) \\
    &= (2\pi\sigma^2)^{-n/2} \exp\left(-\frac{1}{2\sigma^2} \sum_{i=1}^n (y_i - \bm{x}_i^\top \bm{\beta})^2\right)
\end{align*}

\underline{Логарифмическое правдоподобие}.
Удобнее работать с логарифмической функцией правдоподобия:
\begin{align*}
    \ell(\bm{\beta}, \sigma^2) &= \log L(\bm{\beta}, \sigma^2) = -\frac{n}{2} \log(2\pi\sigma^2) - \frac{1}{2\sigma^2} \sum_{i=1}^n (y_i - \bm{x}_i^\top \bm{\beta})^2
\end{align*}

Максимизируем по $\bm{\beta}$ в силу того, что $\bm{\beta}$ входит только во второе слагаемое.
Поскольку $\sigma^2 > 0$, максимизация $\ell(\bm{\beta}, \sigma^2)$ по $\bm{\beta}$ эквивалентна \textbf{минимизации суммы квадратов ошибок}:
\[ \bm{\beta}_{\text{MLE}} = \arg\max_{\bm{\beta}} \ell(\bm{\beta}, \sigma^2) = \arg\min_{\bm{\beta}} \sum_{i=1}^n
(y_i - \bm{x}_i^\top \bm{\beta})^2 \]

\underline{Итог}
\begin{itemize}
    \item При гауссовом предположении о нормальности ошибок \textbf{метод максимального правдоподобия} приводит к
    необходимости минимизировать сумму квадратов ошибок.
    \item Таким образом, минимизация MSE является \textbf{естественным} и \textbf{оптимальным} критерием в рамках
    этой вероятностной модели.
\end{itemize}

\subsection{Аналитическое решение для весов (матричная форма)}

В случае минимизации среднеквадратичной ошибки существует аналитическое решение:
%     (E𝑤p = 𝑤) и оптимальной (то есть имеет наименьшую дисперсию среди всех несмещённых).

\subsubsection{Доказательство}
\begin{gather*}
    \abs{y - X w}^2_2 = (y - X w) (y - X w) = Y^T Y - (X w)^T Y - Y^T(X w) + (X w)^T (X w) = Y^T Y - w^T X^T Y - Y^T X w
    + w^T X^T X w \\
    \nabla_w \abs{y - X w}^2_2 = 0 - X^T Y - X^T Y - (X^T X + X^T X) w
\end{gather*}
Приравнивая $\nabla_w \abs{y - X w}^2_2$ к нулю, получаем $w = (X^T X)^{-1} X^T Y$

Аналитическое решение существует только для 2-нормы.
Для остальных норм решение ищется только градиентными методами

\subsection{Теорема Гаусса–Маркова, BLUE}

\underline{Теорема Гаусса-Маркова}.
Пусть целевые значения выражаются в следующем виде: $Y = X w + \varepsilon$, где $\varepsilon = (\varepsilon_1,
\dots, \varepsilon_N)$ -- случайный вектор, при этом:
\begin{enumerate}
    \item $\mathbb{E} \varepsilon_i = 0 \forall i$
    \item $\mathbb{D} \varepsilon_i = \sigma^2 < \inf \forall i$
    \item $cov(\varepsilon_i, \varepsilon_j) = 0 \forall i \neq j$
\end{enumerate}

Тогда решение задачи наименьших квадратов является оптимальным среди несмещённых.
Другими словами, является несмещённой.

\begin{definition}[BLUE --- Best Linear Unbiased Estimator]
    Оценка $\tilde{\bm{\beta}}$ называется \textbf{наилучшей линейной несмещенной оценкой} (BLUE) вектора параметров
    $\bm{\beta}$, если:
    \begin{enumerate}
        \item \textbf{Линейна:} $\tilde{\bm{\beta}} = \bm{C}\bm{y}$ для некоторой матрицы $\bm{C} \in
        \mathbb{R}^{p \times n}$ (не зависящей от $\bm{y}$)
        \item \textbf{Несмещенна:} $\mathbb{E}[\tilde{\bm{\beta}}] = \bm{\beta}$ для всех $\bm{\beta} \in
        \mathbb{R}^p$
        \item \textbf{Наилучша:} среди всех линейных несмещенных оценок имеет наименьшую дисперсию (в смысле
        положительно определенных матриц ковариаций)
    \end{enumerate}
\end{definition}

Формально, если $\tilde{\bm{\beta}}_1$ и $\tilde{\bm{\beta}}_2$ --- две линейные несмещенные оценки, то
$\tilde{\bm{\beta}}_1$ лучше, если матрица $\text{Cov}(\tilde{\bm{\beta}}_2) - \text{Cov}(\tilde{\bm{\beta}}_1)$
является положительно полуопределенной.